\chapter{Cas expérimental et validation}\label{ch:cas}

\section{Principe des fixtures}

Les fixtures sont des scénarios de test, construits à partir du scénario ZENER et complétés par des
paramètres DDMRP explicites. Elles ne sont pas des paramètres industriels : leur rôle est de faire
apparaître un comportement précis. Chaque fixture isole une situation (stock en zone verte, stock
nul, commande pendant un approvisionnement ouvert, reprise de la non-régression MTS).

\section{Étude de cas principale : DDMRP-C}

La configuration est la suivante : stratégie MTO, une commande client de 10 unités, stocks matière
initialement à zéro, buffer DDMRP actif sur les trois matières. Pour le GPL, le lead time est de
60~h métier et la fenêtre de qualification de la demande est de 400~h. Le délai de première commande
configuré vaut 180~s simulées.

Le déroulement observé dans le run final validé est le suivant. Avant la création de \jv{CMD\_1},
le DDMRP a déjà créé un Open Supply GPL de 4250 unités. Lorsque la commande arrive, on a
$OnHand = 0$, $OpenSupply = 4250$, $QualifiedDemand = 125$, donc
\[
NFP = 0 + 4250 - 125 = 4125.
\]
Cette position est très au-dessus du TOY ; le buffer ne demande donc aucune nouvelle commande. La
réception GPL, de 4250 unités, est physiquement créditée vers $T \approx 556{,}1$. La réservation
sur les trois matières et l'autorisation de Make ont lieu vers $T \approx 713$. La commande est clôturée
vers $T \approx 1876{,}3$, avec l'invariant \jv{[INVARIANT MATIERE] cloture CMD\_1 : OK} et le statut
\Etat{SERVIE} pour 10 unités.

Le résultat essentiel est le suivant : \emph{aucune réception d'origine MTO n'a été créée} pour couvrir
les 125 unités de GPL. La commande a utilisé l'Open Supply DDMRP déjà ouvert. Les temps cités sont des
observations du run final ; les anciennes valeurs de préparation des fixtures ne sont pas reprises.

\begin{figure}[H]
  \centering
  \begin{tikzpicture}[scale=0.9, transform shape, font=\scriptsize, >=Stealth,
    head/.style={draw=ink, rounded corners=2pt, fill=white, align=center, inner sep=2pt, minimum width=1.7cm, minimum height=0.8cm},
    msg/.style={-{Stealth[length=1.8mm]}, thick},
    lab/.style={fill=white, inner sep=0.8pt, font=\scriptsize, align=center}]
    % lignes de vie
    \foreach \xpos/\nomligne/\couleur in {0/Client/ink, 2.1/CommandeAgent/ink, 4.2/MTO/mtoc, 6.3/FicheMatiere/resc,
                          8.4/DDMRP/ddmrpgrn, 10.5/Reception/osc, 12.6/Source/ink, 14.7/Make/mtsc, 16.8/Deliver/ink} {
      \node[head, draw=\couleur] (h\nomligne) at (\xpos,0) {\nomligne};
      \draw[dashed, muted] (\xpos,-0.45) -- (\xpos,-15.6);
    }
    % 1. decision DDMRP
    \node[lab, fill=ddmrpgrn!15] at (8.4,-0.9) {1. NFP inférieure au TOY};
    % 2 à 3. creation de la reception DDMRP
    \draw[msg] (8.4,-1.8) -- (10.5,-1.8) node[lab, midway, above] {2. créer réception DDMRP};
    \draw[msg] (10.5,-2.6) -- (6.3,-2.6) node[lab, midway, above] {3. Open Supply = 4250};
    % 4 à 5. arrivee de la commande
    \draw[msg] (0,-3.4) -- (2.1,-3.4) node[lab, midway, above] {4. CMD\_1 (10 u), T≈348};
    \draw[msg] (2.1,-4.2) -- (4.2,-4.2) node[lab, midway, above] {5. besoin 125 / 10 / 10};
    % 6 a 7. stock insuffisant, consultation Open Supply
    \draw[msg] (4.2,-5.0) -- (6.3,-5.0) node[lab, midway, above] {6. stock physique 0};
    \draw[msg] (4.2,-5.8) -- (10.5,-5.8) node[lab, midway, above] {7. consulter Open Supply existant};
    % 8 a 9. allocation sans reception MTO
    \draw[msg] (10.5,-6.6) -- (4.2,-6.6) node[lab, midway, above] {8. allouer 125 GPL};
    \node[lab, fill=mtoc!12] at (4.2,-7.4) {9. aucune réception MTO créée};
    % 10 a 11. credit physique de la reception DDMRP
    \draw[msg] (10.5,-8.2) -- (6.3,-8.2) node[lab, midway, above] {10. crédit physique T≈556,1 : 4250};
    \node[lab, fill=osc!12] at (6.3,-8.9) {11. Open Supply retiré};
    % 12 a 13. reservation et autorisation Make
    \draw[msg] (4.2,-9.7) -- (6.3,-9.7) node[lab, midway, above] {12. réservation T≈713};
    \draw[msg] (4.2,-10.5) -- (14.7,-10.5) node[lab, midway, above] {13. autorisation Make};
    % 14 a 15. consommation et livraison
    \draw[msg] (14.7,-11.3) -- (6.3,-11.3) node[lab, midway, above] {14. consommation 125 / 10 / 10};
    \draw[msg] (14.7,-12.1) -- (16.8,-12.1) node[lab, midway, above] {15. livraison};
    % 16. invariant final
    \node[lab, fill=resc!15] at (6.3,-13.2) {16. invariant OK, CMD\_1 SERVIE (T≈1876,3)};
  \end{tikzpicture}
  \caption{Séquence complète du cas DDMRP-C, du buffer à la clôture de la commande.
    Les temps sont ceux du run final validé ; les mentions T≈ sont des instants simulés.}
  \label{fig:sequence}
\end{figure}


Une observation mérite d'être notée. La réception GPL est arrivée physiquement vers $T \approx 556{,}1$,
mais la reprise de la commande et sa réservation n'ont eu lieu que vers $T \approx 713$. La vérification
finale de l'approvisionnement avait en effet été planifiée à une échéance ultérieure. Le mécanisme est
correct, puisque la matière est bien réservée avant Make. Il reste toutefois une marge d'optimisation :
une reprise événementielle, lancée dès la réception d'un Open Supply alloué, éviterait cette attente.
Cette piste est traitée au \cref{ch:apports}.

Les instants de cette chronologie sont de deux natures. Le run a démarré effectivement à $T \approx 166{,}1$
sur l'horloge absolue d'AnyLogic. Le délai de première commande, $T+181{,}9$, est relatif à ce démarrage
effectif. La création de \jv{CMD\_1} est donc observée à l'instant absolu
\[
166{,}1 + 181{,}9 \approx 348{,}0,
\]
ce qui correspond au $T \approx 348$ des observations. Cette distinction entre temps relatif et temps
absolu est celle de la section « Temps métier et temps de simulation » du \cref{ch:ddmrp}.

\section{Synthèse de la campagne}

\begin{table}[H]
  \centering
  \caption{Résultats de la campagne runtime validée.}\label{tab:campagne}
  \begin{tabularx}{\textwidth}{@{}L{3.2cm}L{2.4cm}Y@{}}
    \toprule
    Essai & Résultat & Observation \\
    \midrule
    MTO-2 & PASS & réservation et consommation avec invariant \\
    DDMRP-A & PASS & stock initial en zone verte : aucune commande \\
    DDMRP-B & PASS & stock nul : création d'un approvisionnement, puis retour en zone verte ;
      GPL 500, bouteilles 80 (TOG $\approx$ 71{,}73 arrondi au MOQ 10), accessoires 40 \\
    DDMRP-C & PASS complet & MTO utilisant un Open Supply DDMRP préexistant, sans double approvisionnement \\
    Non-régression MTS & PASS fonctionnel & REAPPRO\_1 produit 10/10 ; invariants OK ; CMD\_1 livrée et clôturée \\
    \bottomrule
  \end{tabularx}
\end{table}

Dans le test de non-régression, la commande a été livrée en retard. Ce test est volontairement sévère :
le stock de produits finis et les matières sont à zéro, et la livraison attend la reconstitution complète du
stock. Ce retard résulte donc de la reconstruction du stock, et ne signale pas une panne du mécanisme. Aucune
logique DDMRP n'était active dans ce run.

Une anomalie a été corrigée pendant la campagne : les réceptions des ordres REAPPRO étaient étiquetées
\texttt{origine=MTO} alors qu'elles relèvent de commandes MTS. L'étiquette reflète désormais le type réel de
la commande. Seule la trace est modifiée : quantités, allocations, échéances et décisions sont inchangées.

\section{Coexistence avec la politique (s,Q)}

Une matière est pilotée soit par DDMRP, soit par la politique (s,Q) historique, jamais par les deux. Dans
\jv{calculerBesoinsNets()}, les réceptions échues sont d'abord traitées. Si la matière est gérée par DDMRP,
la décision DDMRP est prise puis la boucle passe à la matière suivante. Sinon, la politique (s,Q) historique
s'applique (\cref{fig:arbre}).

\begin{figure}[H]
  \centering
  \begin{tikzpicture}[font=\small, >=Stealth, node distance=1.0cm,
    dec/.style={draw=ink, diamond, aspect=2.4, align=center, inner sep=2pt, fill=white},
    act/.style={draw=ink, rounded corners=3pt, align=center, inner sep=3pt, fill=white, text width=3.0cm},
    tr/.style={-{Stealth[length=2mm]}, thick}]
    \node[act] (rcv) {Traiter les réceptions\\échues de la matière};
    \node[dec, below=of rcv] (ddm) {Matière gérée\\par DDMRP ?};
    \node[act, fill=ddmrpgrn!20, below=of ddm, xshift=-3.8cm] (dec) {\texttt{decisionDDMRP()}\\NFP, zones, ordre DDMRP};
    \node[act, below=1.2cm of dec] (end1) {Sortie : (s,Q) non appliquée\\à cette matière};
    \node[dec, below=of ddm, xshift=3.8cm] (rea) {Production autonome\\REAPPRO en cours ?};
    \node[act, below=1.2cm of rea, fill=ink!6] (skip) {Passer la décision (s,Q)\\pour cette matière};
    \node[act, fill=mtsc!15, right=0.9cm of rea] (sq) {Politique (s,Q) historique\\point de commande, Q*};
    \draw[tr] (rcv) -- (ddm);
    \draw[tr] (ddm) -- node[above, font=\scriptsize] {oui} (dec);
    \draw[tr] (dec) -- (end1);
    \draw[tr] (ddm) -- node[above, font=\scriptsize] {non} (rea);
    \draw[tr] (rea) -- node[above, font=\scriptsize] {oui} (skip);
    \draw[tr] (rea) -- node[above, font=\scriptsize] {non} (sq);
  \end{tikzpicture}
  \caption{Coexistence : une matière est pilotée soit par DDMRP, soit par la politique (s,Q),
    jamais par les deux.}
  \label{fig:arbre}
\end{figure}


\section{Observabilité}

Le système se rend observable par trois canaux. Les journaux comportent les étiquettes \texttt{[DDMRP]} (une
ligne par décision, avec NFP, zones et action), \texttt{[OPEN SUPPLY]} (création et crédit de réceptions),
\texttt{[MATIERE]} (réservations) et \texttt{[INVARIANT MATIERE]} (contrôles à la clôture). La feuille
\texttt{DDMRP} du tableur d'export donne, pour chaque matière, le stock physique, les réservations, la
disponibilité opérationnelle, l'Open Supply, la demande qualifiée et ses sous-totaux, la NFP, l'ADU, le lead
time, les zones, les seuils TOR, TOY et TOG, la quantité recommandée, le statut et la priorité. Les mêmes
grandeurs sont exportées dans le graphe ABox sous la forme de triplets \texttt{run:ddmrp*}, ce qui permet
une exploitation sémantique ultérieure.
